\newcommand{\DoNotLoadEpstopdf}{}
\documentclass[11pt,reqno]{amsart}
\usepackage[T1]{fontenc}
\usepackage{lmodern,amsmath,amssymb,amsthm,mathtools,mathrsfs,microtype}
\usepackage[margin=1.1in]{geometry}
\usepackage{xr-hyper}
\usepackage[hidelinks]{hyperref}
\externaldocument[supp-]{build/supplement}[supplement.pdf]
\hypersetup{pdftitle={Collision records, arithmetic, and raw local inversion in a triangular Lorentz gas},pdfauthor={Qian Qi}}
\setlength{\emergencystretch}{2em}
\allowdisplaybreaks[2]
\numberwithin{equation}{section}
\newtheorem{theorem}{Theorem}[section]
\newtheorem{lemma}[theorem]{Lemma}
\newtheorem{proposition}[theorem]{Proposition}
\newtheorem{corollary}[theorem]{Corollary}
\theoremstyle{definition}
\newtheorem{definition}[theorem]{Definition}
\newtheorem{inputstatement}[theorem]{Input statement}
\theoremstyle{remark}\newtheorem{remark}[theorem]{Remark}
\newcommand{\R}{\mathbb R}\newcommand{\Z}{\mathbb Z}
\newcommand{\one}{\mathbf1}\newcommand{\dd}{\,\mathrm d}
\newcommand{\esssup}{\operatorname*{ess\,sup}}
\newcommand{\diver}{\operatorname{div}}\newcommand{\TV}{\operatorname{TV}}
\title[Collision records and raw local inversion]{Collision records, arithmetic, and raw local inversion in a triangular Lorentz gas}
\author{Qian Qi}
\date{October 13, 2026. A2-DYN, revision 102}
\subjclass[2020]{37D50, 37A50, 60F05, 28A75}
\keywords{Lorentz gas, exact collision records, paired physical traces, transverse coarea, pointwise inversion}
\begin{document}
\raggedbottom
\begin{abstract}
We study the exact-label pointwise raw local limit problem for actual
returns in a triangular Lorentz gas. The remaining upper source is
the positive minimum of two physical endpoint traces on each short
regular segment. We represent this minimum by a barycentric coarea
measure and construct a transverse comparison which preserves relative
position between the two ends. Its Jacobian is computed before
physical sheets are summed. A capped transverse clock then gives an
essential-height estimate for the admitted positive paired source
itself, rather than for its outgoing-minus-incoming difference.
A relative endpoint certificate applies without a lower residence
cutoff. Under the original reconstruction schedule the admitted
contribution is $O(B^{-1/384})$. The unadmitted transverse source is
retained explicitly; no assertion that it vanishes on the full Lorentz
family is made. The original pointwise objective, exact labels,
physical guard, arithmetic classes and all historical proofs are
unchanged.
\end{abstract}
\maketitle
\input{journal-problem-v93}

\subsection{Positive paired mass and a transverse comparison}
The uniform momentum-capacity theorem has removed the singular
volume-source premise from the preceding scalar reduction. Its
remaining boundary quantity is positive: two large endpoint traces
can have a small difference. We therefore put the positive minimum
itself into a coarea measure. The new comparison moves in angular
direction, keeping fixed the relative position between the same two
physical ends. It does not require a radial trajectory to survive
beyond its original short residence.

Theorem~\ref{thm:v102-pair-coarea} constructs this measure and its
Jacobian. Theorem~\ref{thm:v102-admitted-height} estimates its whole
admitted part in essential height after physical summation, with no
word-count or chart-overlap factor. Theorem~\ref{thm:v102-certificate}
gives a test in terms of the two endpoint roof maps and their positive
trace density. Its constants do not deteriorate merely because the
original residence tends to zero. Theorem~\ref{thm:v102-scalar}
inserts the estimate into the unchanged scalar target.

Admission on the complete nongrazing paired family is not established.
The remaining positive source includes short transverse comparison
lifetimes, folds of the endpoint roof maps, and insufficient control
of relative trace variation. A backward critical trace is not set to
zero. These distinctions are part of the theorem statements, not
changes of the governing objective.

\input{journal-inputs-v102}
\input{core/244_barycentric_positive_trace_coarea}
\input{core/245_transverse_clock_and_pointwise_trace_height}
\input{core/246_transverse_certificates_and_scalar_schedule}

\appendix
\input{appendices/v102_dependency_map}
\input{journal-references-v100}
\end{document}
